% GENERATED FILE. Edit canonical lessons, then run npm run generate:books.
% canonical-source-hash: fnv1a64:d2cae34ef528c0ef
% canonical-lessons: ML-C499-vanam, ML-C499-vanyajivi, ML-C499-sanketam, ML-C499-kattana, ML-C499-kattupottu

\chapter{In the Forest}
\label{ch:ml-in-the-forest}
\hlchaptermodality{499}

\begin{chapteropening}
\textbf{By the end of this chapter:} \emph{I can say a forest, a wild animal, a sanctuary, a wild elephant and a gaur.}

Complete the last lesson of chapter 499: I can say a forest, a wild animal, a sanctuary, a wild elephant and a gaur.
\end{chapteropening}

\section[\texorpdfstring{\ml{വനം} (vanaṁ)}{വനം (vanaṁ)}]{\ml{വനം} (vanaṁ) — a forest}

\label{lesson:ML-C499-vanam}



\begin{quote}
Before the new one: say the Malayalam for a shark, then the Malayalam for to accumulate.
\end{quote}

\subsection*{You'll want to know: \ml{വനം}}
\textbf{\ml{വനം}} — \emph{vanaṁ} — \textquotedblleft{}a forest\textquotedblright{}.

\begin{quote}
\ml{ആ} \ml{വനത്തിൽ} \ml{ധാരാളം} \ml{ആനകൾ} \ml{ഉണ്ട്}.
\end{quote}

\begin{quote}
\emph{ā vanattil dhārāḷaṁ ānakaḷ uṇṭŭ.}
\end{quote}

\begin{quote}
There are many elephants in that forest.
\end{quote}

A formal word; in speech people say \ml{കാട്}. The forest department is the \ml{വനംവകുപ്പ്}.

\begin{cousinweb}
\textbf{\ml{വനം}} (\emph{vanaṁ}) is borrowed from Sanskrit \emph{vana}, \textquotedblleft{}a forest, a wood\textquotedblright{} (Monier-Williams).
\end{cousinweb}

\begin{grammarlens}[title={What you've built}]
One more word for getting around.
\end{grammarlens}

\subsection*{Your turn}
\begin{itemize}
\raggedright
  \item \emph{Say it:} \emph{vanaṁ}
  \item \emph{Say it:} \emph{vanaṁ}, once more
  \item \emph{Say it:} \emph{ā vanattil dhārāḷaṁ ānakaḷ uṇṭŭ.} — There are many elephants in that forest.
  \item \emph{Recall:} say the Malayalam for a shark, then the Malayalam for to accumulate, then say \emph{vanaṁ} again
\end{itemize}

\begin{tcolorbox}[breakable,colback=teal!4,colframe=teal!35!black,title={Before you move on}]
What does \ml{വനം} mean? (\textquotedblleft{}A forest\textquotedblright{}.) Say it once more.
\end{tcolorbox}
\section[\texorpdfstring{\ml{വന്യജീവി} (vanyajīvi)}{വന്യജീവി (vanyajīvi)}]{\ml{വന്യജീവി} (vanyajīvi) — a wild animal; wildlife}

\label{lesson:ML-C499-vanyajivi}



\begin{quote}
Before the new one: say the Malayalam for to accumulate, then the Malayalam for a forest.
\end{quote}

\subsection*{You'll want to know: \ml{വന്യജീവി}}
\textbf{\ml{വന്യജീവി}} — \emph{vanyajīvi} — \textquotedblleft{}a wild animal; wildlife\textquotedblright{}.

\begin{quote}
\ml{വനത്തിൽ} \ml{വന്യജീവികളെ} \ml{കാണാൻ} \ml{ഞങ്ങൾ} \ml{രാവിലെ} \ml{പോയി}.
\end{quote}

\begin{quote}
\emph{vanattil vanyajīvikaḷe kāṇān ñaṅṅaḷ rāvile pōyi.}
\end{quote}

\begin{quote}
We went in the morning to see wild animals in the forest.
\end{quote}

Plural: \ml{വന്യജീവികൾ}, \emph{vanyajīvikaḷ}.

\ml{വന്യ}, wild, of the forest, plus \ml{ജീവി}, a living thing.

\begin{cousinweb}
\textbf{\ml{വന്യജീവി}} (\emph{vanyajīvi}) is a Malayalam compound built from two Sanskrit words, \emph{vanya}, \textquotedblleft{}growing or living in a forest, wild\textquotedblright{}, and \emph{jīvin}, \textquotedblleft{}living\textquotedblright{} (Monier-Williams).
\end{cousinweb}

\begin{grammarlens}[title={What you've built}]
One more word for getting around.
\end{grammarlens}

\subsection*{Your turn}
\begin{itemize}
\raggedright
  \item \emph{Say it:} \emph{vanyajīvi}
  \item \emph{Say it:} \emph{vanyajīvi}, once more
  \item \emph{Say it:} \emph{vanattil vanyajīvikaḷe kāṇān ñaṅṅaḷ rāvile pōyi.} — We went in the morning to see wild animals in the forest.
  \item \emph{Recall:} say the Malayalam for to accumulate, then the Malayalam for a forest, then say \emph{vanyajīvi} again
\end{itemize}

\begin{tcolorbox}[breakable,colback=teal!4,colframe=teal!35!black,title={Before you move on}]
What does \ml{വന്യജീവി} mean? (\textquotedblleft{}A wild animal; wildlife\textquotedblright{}.) Say it once more.
\end{tcolorbox}
\section[\texorpdfstring{\ml{സങ്കേതം} (saṅkētaṁ)}{സങ്കേതം (saṅkētaṁ)}]{\ml{സങ്കേതം} (saṅkētaṁ) — a sanctuary; a refuge}

\label{lesson:ML-C499-sanketam}



\begin{quote}
Before the new one: say the Malayalam for a forest, then the Malayalam for a wild animal.
\end{quote}

\subsection*{You'll want to know: \ml{സങ്കേതം}}
\textbf{\ml{സങ്കേതം}} — \emph{saṅkētaṁ} — \textquotedblleft{}a sanctuary; a refuge\textquotedblright{}.

\begin{quote}
\ml{ആ} \ml{പക്ഷിസങ്കേതത്തിൽ} \ml{പല} \ml{തരം} \ml{പക്ഷികൾ} \ml{ഉണ്ട്}.
\end{quote}

\begin{quote}
\emph{ā pakṣisaṅkētattil pala taraṁ pakṣikaḷ uṇṭŭ.}
\end{quote}

\begin{quote}
That bird sanctuary has many kinds of birds.
\end{quote}

A protected place: \ml{പക്ഷിസങ്കേതം}, a bird sanctuary; \ml{വന്യജീവിസങ്കേതം}, a wildlife sanctuary.

\begin{cousinweb}
\textbf{\ml{സങ്കേതം}} (\emph{saṅkētaṁ}) is borrowed from Sanskrit \emph{saṅketa}, \textquotedblleft{}an agreed sign; an appointed place\textquotedblright{} (Monier-Williams); in Malayalam it came to mean a place set aside.
\end{cousinweb}

\begin{grammarlens}[title={What you've built}]
One more word for getting around.
\end{grammarlens}

\subsection*{Your turn}
\begin{itemize}
\raggedright
  \item \emph{Say it:} \emph{saṅkētaṁ}
  \item \emph{Say it:} \emph{saṅkētaṁ}, once more
  \item \emph{Say it:} \emph{ā pakṣisaṅkētattil pala taraṁ pakṣikaḷ uṇṭŭ.} — That bird sanctuary has many kinds of birds.
  \item \emph{Recall:} say the Malayalam for a forest, then the Malayalam for a wild animal, then say \emph{saṅkētaṁ} again
\end{itemize}

\begin{tcolorbox}[breakable,colback=teal!4,colframe=teal!35!black,title={Before you move on}]
What does \ml{സങ്കേതം} mean? (\textquotedblleft{}A sanctuary; a refuge\textquotedblright{}.) Say it once more.
\end{tcolorbox}
\section[\texorpdfstring{\ml{കാട്ടാന} (kāṭṭāna)}{കാട്ടാന (kāṭṭāna)}]{\ml{കാട്ടാന} (kāṭṭāna) — a wild elephant}

\label{lesson:ML-C499-kattana}



\begin{quote}
Before the new one: say the Malayalam for a wild animal, then the Malayalam for a sanctuary.
\end{quote}

\subsection*{You'll want to know: \ml{കാട്ടാന}}
\textbf{\ml{കാട്ടാന}} — \emph{kāṭṭāna} — \textquotedblleft{}a wild elephant\textquotedblright{}.

\begin{quote}
\ml{രാത്രി} \ml{ഒരു} \ml{കാട്ടാന} \ml{വയലിൽ} \ml{ഇറങ്ങി}.
\end{quote}

\begin{quote}
\emph{rātri oru kāṭṭāna vayalil iṟaṅṅi.}
\end{quote}

\begin{quote}
At night a wild elephant came down into the fields.
\end{quote}

Plural: \ml{കാട്ടാനകൾ}, \emph{kāṭṭānakaḷ}.

\ml{കാട്}, the forest, plus \ml{ആന}: a forest elephant, not a temple one.

\begin{cousinweb}
\textbf{\ml{കാട്ടാന}} (\emph{kāṭṭāna}) is a native Dravidian compound of \ml{കാട്}, a forest (DEDR 1438), and \ml{ആന}, an elephant (DEDR 5161).
\end{cousinweb}

\begin{grammarlens}[title={What you've built}]
One more word for getting around.
\end{grammarlens}

\subsection*{Your turn}
\begin{itemize}
\raggedright
  \item \emph{Say it:} \emph{kāṭṭāna}
  \item \emph{Say it:} \emph{kāṭṭāna}, once more
  \item \emph{Say it:} \emph{rātri oru kāṭṭāna vayalil iṟaṅṅi.} — At night a wild elephant came down into the fields.
  \item \emph{Recall:} say the Malayalam for a wild animal, then the Malayalam for a sanctuary, then say \emph{kāṭṭāna} again
\end{itemize}

\begin{tcolorbox}[breakable,colback=teal!4,colframe=teal!35!black,title={Before you move on}]
What does \ml{കാട്ടാന} mean? (\textquotedblleft{}A wild elephant\textquotedblright{}.) Say it once more.
\end{tcolorbox}
\section[\texorpdfstring{\ml{കാട്ടുപോത്ത്} (kāṭṭupōttŭ)}{കാട്ടുപോത്ത് (kāṭṭupōttŭ)}]{\ml{കാട്ടുപോത്ത്} (kāṭṭupōttŭ) — a gaur; an Indian bison}

\label{lesson:ML-C499-kattupottu}



\begin{quote}
Before the new one: say the Malayalam for a sanctuary, then the Malayalam for a wild elephant.
\end{quote}

\subsection*{You'll want to know: \ml{കാട്ടുപോത്ത്}}
\textbf{\ml{കാട്ടുപോത്ത്}} — \emph{kāṭṭupōttŭ} — \textquotedblleft{}a gaur; an Indian bison\textquotedblright{}.

\begin{quote}
\ml{കാട്ടുപോത്ത്} \ml{ഒരു} \ml{വലിയ} \ml{കാട്ടുമൃഗമാണ്}.
\end{quote}

\begin{quote}
\emph{kāṭṭupōttŭ oru valiya kāṭṭumṛgamāṇŭ.}
\end{quote}

\begin{quote}
The gaur is a big wild animal.
\end{quote}

\ml{കാട്ടു}-, of the forest, plus \ml{പോത്ത്}, a buffalo.

\begin{cousinweb}
\textbf{\ml{കാട്ടുപോത്ത്}} (\emph{kāṭṭupōttŭ}) is a native Dravidian compound of \ml{കാട്}, a forest (DEDR 1438), and \ml{പോത്ത്}, a male buffalo (DEDR 4586).
\end{cousinweb}

\begin{grammarlens}[title={What you've built}]
One more word for getting around.
\end{grammarlens}

\subsection*{Your turn}
\begin{itemize}
\raggedright
  \item \emph{Say it:} \emph{kāṭṭupōttŭ}
  \item \emph{Say it:} \emph{kāṭṭupōttŭ}, once more
  \item \emph{Say it:} \emph{kāṭṭupōttŭ oru valiya kāṭṭumṛgamāṇŭ.} — The gaur is a big wild animal.
  \item \emph{Recall:} say the Malayalam for a sanctuary, then the Malayalam for a wild elephant, then say \emph{kāṭṭupōttŭ} again
\end{itemize}

\begin{tcolorbox}[breakable,colback=teal!4,colframe=teal!35!black,title={Before you move on}]
What does \ml{കാട്ടുപോത്ത്} mean? (\textquotedblleft{}A gaur; an Indian bison\textquotedblright{}.) Say it once more.
\end{tcolorbox}
